\chapter{Conclusion and Discussion}
\label{chap:conclusion-discussion}

\noindent
This chapter concludes the report by reflecting on the project outcomes, summarizing key challenges encountered during development and evaluation, and outlining directions for future improvement.
\vspace{1em}

\section{Reflection}
This project set out to answer a constrained research question---whether prestimulus EEG alone can predict behavioral performance---and to deliver software that makes this question testable in a reproducible way. The main contribution is an end-to-end workflow that links (1) event-locked dataset construction for CCD-ITI windows, (2) standardized preprocessing and feature extraction for steady-state characteristics, (3) dual-target modeling (RT regression and correctness classification), and (4) consistent evaluation and reporting artifacts.

From a software engineering perspective, the outcome is a clearer separation between data ingestion, processing, modeling, and reporting. This reduces the risk that results depend on hidden notebook state and makes it easier to rerun experiments, compare configurations, and regenerate figures/tables for the final report.

From a research perspective, the journey was difficult: substantial effort was required before model training became stable and comparable. Even after overcoming pipeline and data-quality challenges, the final models still did not achieve strong predictive performance.

\section{Challenges}
The main challenges encountered are:
\begin{itemize}
    \item \textbf{Trial and label alignment:} prestimulus prediction depends on strict alignment between events, ITI windows, and behavioral labels; small mistakes can invalidate comparisons.
    \item \textbf{Data quality and non-stationarity:} EEG recordings vary across subjects and sessions, and artifacts/noise can dominate unless preprocessing and sanity checks are systematic.
    \item \textbf{Class imbalance and instability for correctness:} several runs showed strong sensitivity but weak specificity, indicating prediction collapse toward one class despite acceptable overall accuracy.
    \item \textbf{Imbalance mitigation with limited gain:} class weighting, weighted sampling, and undersampling were tested to correct imbalance effects, but improvements were partial and did not deliver consistently reliable classification quality.
    \item \textbf{Weak RT explanatory power:} regression runs with MAE/MSE/Huber losses all produced near-zero R², suggesting underfitting or insufficient predictive signal extraction in the current setup.
    \item \textbf{Reproducibility overhead:} making experiments repeatable requires extra engineering (configs, caching, logging, deterministic splits) beyond what is needed for a one-off notebook result.
\end{itemize}

\section{Future Work}
Future work should prioritize making the platform easier to extend and easier to use in real settings. Specific next steps address the observed limitations:

\subsection*{Platform Usability and Extensibility (Primary Priority)}
\begin{itemize}
    \item \textbf{Plugin-ready architecture:} Define clear interfaces for preprocessing, feature extraction, and modeling so new steps can be added without modifying core code. Provide reference implementations and documentation for each extension point.
    \item \textbf{User onboarding and documentation:} Write step-by-step tutorials, a user manual with screenshots, and documented example workflows for common tasks (loading BIDS data, running preprocessing, launching a model). Include troubleshooting guides.
    \item \textbf{Usability testing and feedback loops:} Conduct user testing sessions with domain researchers to identify confusing UI steps. Collect satisfaction surveys and iterate on the interface based on real feedback.
    \item \textbf{Operational clarity:} Improve error messages with actionable guidance, add run summaries showing what succeeded and what failed, and keep structured experiment logs (inputs, outputs, parameters, timing) for reproducible comparisons and failure auditing.
\end{itemize}

\subsection*{Model-Focused Improvements (Secondary Priority)}
\begin{itemize}
    \item \textbf{Fix class imbalance and evaluation metrics:} The correctness classifier showed 95\% sensitivity but only \textasciitilde 10\% specificity on error detection. Replace accuracy with balanced metrics (F1, MCC) during training and hyperparameter selection. Test threshold tuning and probability calibration post-hoc.
    \item \textbf{Improve RT regression capacity:} Current MAE/MSE/Huber experiments all produced R\textsuperscript{2} near zero. Explore richer feature representations (time-frequency, connectivity) and stronger regularization schedules to move beyond mean-prediction behavior.
\end{itemize}
